\documentclass{article}

\usepackage[margin=0.5in]{geometry}
\usepackage{amsmath, mathtools, array, booktabs, hyperref}

\title{HW 3}

\author{Darin Peries}

\date{\today}

\begin{document}
\maketitle

\section*{Problem 1}

\subsection*{Part 1}

\noindent Since $A$ is a square matrix, its characteristic polynomial is defined as follows.
$$p_A(\lambda)=\det(A-I\lambda)$$

\noindent The term $A-I\lambda$ is supposed to result in a singular matrix, so the determinant in the equation above
is supposed to equal zero. The diagonals of $A-I\lambda$ are in the form $a_{ii}-\lambda$. \\

\noindent Because $A$ is a square matrix, its eigenvalues are the diagonals of its upper triangular
matrix factor, and the determinant of an upper triangular matrix, obtaind through gaussian elimination (GE), the product of its diagonal. The equation is as 
follows.

$$\det(GE(A-I\lambda))=(u_{i,i} - \lambda)(u_{i+1,i+1}-\lambda)...(u_{n,n}-\lambda)$$

\noindent To ensure that the determinant is non-zero, set $\lambda=0$. Additionally, since the diagonals of the upper
triangular matrix are the eigen values of $A$, $u_{i,i} = \lambda_i$. The resulting equation is:

$$p_A(0)=\det(GE(A-I(0)))=(\lambda_i)(\lambda_{i+1})...(\lambda_n)=\prod_{i=1}^{n}\lambda_i$$

\subsection*{Part 2}
The iteration matrix is defined as $B_w = (D+wL)^{-1}[(1-w)D-wU]$. By taking the determinant of $B_w$,
the following expression is formed.

$$\det(B_w) = \det(D+wL)^{-1}[\det((w-1)D-wU)] = \det(D)^{-1}[\det((1-w)D)]$$

\noindent Using part 1, the determinant of $D$ is the product of its diagonal entries or eigenvalues.

$$\det(B_w) = \det(D)^{-1}[\det((1-w)D)] = \prod_{i=1}^{n} \frac{1}{\lambda_i} \prod_{i=1}^{n} (1-w)\lambda_i = \prod_{i=1}^{n}\frac{(w-1)\lambda_i}{\lambda_i} = \prod_{i=1}^{n} (1-w)$$

\noindent Since the determinant is the product of a square matrices eigenvalues, the eigenvalues of are $\Lambda(B_w)=(1-w)$.
The spectral radius, $\rho(B_w)$ is bounded below by the gemoetric mean of the eigen values, $\sqrt[n]{\prod_{i=1}^{n}|\lambda_i|}$,
so $\rho(B_w) \ge \sqrt[n]{|(w-1)|^n}$; or more concisely as $\mathbf{\rho(B_w) \ge |w-1|}$.

\section*{Problem 2}

\subsection*{Part 1}

$$B_J = \begin{pmatrix}
    0 & -1/3 & 0 \\ -1/3 & 0 & -1/3 \\ 0 & -1/3 & 0\\
\end{pmatrix}$$

$$B_{GS} = \begin{pmatrix} 
    0 & -1/3 & 0 \\ 0 & 1/9 & -1/3 \\ 0 & -1/27 & 1/9 \\
\end{pmatrix}$$

\noindent The infinity norm of $B_J$ is $\|B_J\|_{\infty} = \max(|-1/3|, |-1/3|+|-1/3|, |-1/3|)=2/3$. $\rho(B_J) = \max\limits_\lambda(\lambda (\lambda^2-2/9)=0)=\frac{\sqrt{2}}{3}$.\\
\noindent The infinity norm of $B_{GS}$ is $\|B_{GS}\|_\infty = \max(|-1/3|, |1/9|+|-1/3|,|-1/27|+|1/9|)=4/9$. $\rho(B_{GS})=\max\limits_\lambda (\lambda(\lambda^2+2/9\lambda+2/81)=0)=2/9$.\\

\noindent The optimal $\omega$ for SOR is $\frac{2}{1+\sqrt{1-\rho(B_J)^2}}=9-3\sqrt{7}$. This only works because A is tridiagonal and strictly diagonally dominant, ie.) symmetric positive definite.\\

\subsection*{Part 2}

\subsubsection*{Tables}

\begin{table}[htbp]
\centering
\caption{Jacobi Method Iterations}
\label{tab:jacobi}
\begin{tabular}{|c c c c c c|}
\toprule
$k$ & $x_1$ & $x_2$ & $x_3$ & $e$ & $e_{\text{suc}}$ \\
\midrule
1  & 1.333333 & 1.666667 & 1.333333 & 0.666667 & 0.000000 \\
2  & 0.777778 & 0.777778 & 0.777778 & 0.222222 & 0.333333 \\
3  & 1.074074 & 1.148148 & 1.074074 & 0.148148 & 0.666667 \\
4  & 0.950617 & 0.950617 & 0.950617 & 0.049383 & 0.333333 \\
5  & 1.016461 & 1.032922 & 1.016461 & 0.032922 & 0.666667 \\
6  & 0.989026 & 0.989026 & 0.989026 & 0.010974 & 0.333333 \\
7  & 1.003658 & 1.007316 & 1.003658 & 0.007316 & 0.666667 \\
8  & 0.997561 & 0.997561 & 0.997561 & 0.002439 & 0.333333 \\
9  & 1.000813 & 1.001626 & 1.000813 & 0.001626 & 0.666667 \\
10 & 0.999458 & 0.999458 & 0.999458 & 0.000542 & 0.333333 \\
\bottomrule
\end{tabular}
\end{table}

\begin{table}[htbp]
\centering
\caption{Gauss--Seidel Method Iterations}
\label{tab:gs}
\begin{tabular}{|c c c c c c|}
\toprule
$k$ & $x_1$ & $x_2$ & $x_3$ & $e$ & $e_{\text{suc}}$ \\
\midrule
1  & 1.333333 & 1.222222 & 0.925926 & 0.333333 & 0.000000 \\
2  & 0.925926 & 1.049383 & 0.983539 & 0.074074 & 0.222222 \\
3  & 0.983539 & 1.010974 & 0.996342 & 0.016461 & 0.222222 \\
4  & 0.996342 & 1.002439 & 0.999187 & 0.003658 & 0.222222 \\
5  & 0.999187 & 1.000542 & 0.999819 & 0.000813 & 0.222222 \\
6  & 0.999819 & 1.000120 & 0.999960 & 0.000181 & 0.222222 \\
7  & 0.999960 & 1.000027 & 0.999991 & 0.000040 & 0.222222 \\
8  & 0.999991 & 1.000006 & 0.999998 & 0.000009 & 0.222222 \\
9  & 0.999998 & 1.000001 & 1.000000 & 0.000002 & 0.222222 \\
10 & 1.000000 & 1.000000 & 1.000000 & 0.000000 & 0.222222 \\
\bottomrule
\end{tabular}
\end{table}


\begin{table}[htbp]
\centering
\caption{SOR Method Iterations ($\omega = 9 - 3\sqrt{7} \approx 1.1707$)}
\label{tab:sor}
\begin{tabular}{|c c c c c c|}
\toprule
$k$ & $x_1$ & $x_2$ & $x_3$ & $e$ & $e_{\text{suc}}$ \\
\midrule
1  & 1.416995 & 1.269275 & 0.967356 & 0.416995 & 0.000000 \\
2  & 0.878445 & 1.037729 & 0.988683 & 0.121555 & 0.291503 \\
3  & 0.994262 & 1.003675 & 0.999408 & 0.005738 & 0.047208 \\
4  & 0.999058 & 1.000313 & 0.999926 & 0.000942 & 0.164097 \\
5  & 0.999948 & 1.000025 & 0.999996 & 0.000052 & 0.054849 \\
6  & 0.999994 & 1.000002 & 1.000000 & 0.000006 & 0.107079 \\
7  & 1.000000 & 1.000000 & 1.000000 & 0.000000 & 0.057452 \\
8  & 1.000000 & 1.000000 & 1.000000 & 0.000000 & 0.091118 \\
9  & 1.000000 & 1.000000 & 1.000000 & 0.000000 & 0.058765 \\
10 & 1.000000 & 1.000000 & 1.000000 & 0.000000 & 0.083607 \\
\bottomrule
\end{tabular}
\end{table}

\subsubsection*{Discussion}
The jacobi method, while it minimizes the error, does not get the lowest possible error. Gauss-Seidel reaches the lowest error at the last 
iteration. For the SOR method, the solution converge smuch earlier than the Gauss-Seidel method.

\section*{Problem 3}

\subsection*{Part 1}

$$x^{(k+1)} = Q^{-1}(Q-A)x^{(k)}+Q^{-1}b \rightarrow x^{(k+1)} = x^{(k)}-AQ^{-1}x^{(k)}+Q^{-1}b $$

$$z^{(k)} = Q^{-1}(b-Ax^{(k)})$$

$$x^{(k+1)}=x^{(k)} + Q^{-1}(b-Ax^{(k)}) \rightarrow x^{(k+1)} = x^{(k)} + z^{(k)}$$

\noindent From the equations above, the iteration process is equuivalent.

\subsection*{Part 2}

\subsubsection*{proof for z}

\[z^{(k+1)} = Q^{-1}r^{(k+1)} = Q^{-1}(b-Ax^{(k+1)}) = Q^{-1}(b-A(x^{(k)}+z^{(k)})) =Q^{-1}(b-Ax^{(k)} - Az^{(k)}) =\] 
\[Q^{-1}(r^{(k)}-Az^{(k)}) = Q^{-1}(Qz^{(k)}-Az^{(k)})=(I-Q^{-1}A)z^{(k)} \]

\subsubsection*{proof for r}
\[r^{(k+1)}=b-Ax^{(k+1)}=b-A(x^{(k)}+z^{(k)})=b-Ax^{(k)}-Az^{(k)}=r^{(k)}-AQ^{-1}r^{(k)}=(I-AQ^{-1})r^{(k)}\]

\section*{Problem 4}

\begin{table}[htbp]
\centering
\caption{Comparison of Jacobi and Gauss--Seidel Convergence}
\label{tab:comparison}
\begin{tabular}{|l c c c|}
\toprule
Method & $k$ & $r_{\text{norm}}$ & $r_{\text{successive}}$ \\
\midrule
Jacobi        & 9 & 0.007820 & 0.500488 \\
Gauss--Seidel & 5 & 0.009747 & 0.502953 \\
\bottomrule
\end{tabular}
\end{table}

\subsubsection*{Discussion}
The results of the table show that the Gauss-Siedel method converges much faster than the Jacobi. This is due to the fact that the 
spectral radius of the Gauss-Seidel matrix is the spectral radius of the jacobi matrix squared. Since both spectral radii are less than one,
the Gauss-Seidel method converges much faster. Therefore, the behavior in the table above matches what we've learned in class.

\end{document}